
This work intersects three research areas: AI alignment methods, user adaptation in human-computer interaction, and conversation-level quality metrics.

\subsubsection{AI Alignment and RLHF}

Reinforcement Learning from Human Feedback has become the dominant paradigm for aligning large language models with human preferences (Christiano et al., 2017; Ouyang et al., 2022). The standard approach trains a reward model on human preference judgments, then uses reinforcement learning to optimize the language model policy to maximize predicted reward. InstructGPT \citep{ouyang2022training} demonstrated that RLHF-trained models produce outputs humans rate as more helpful, honest, and harmless. Constitutional AI \citep{bai2022constitutional} extended this framework by incorporating AI-generated feedback.

These methods treat alignment as unidirectional optimization: adjust AI policy until human evaluators are satisfied. Evaluation focuses on aggregate preference accuracy or single-turn response quality. Success is measured by how well the AI conforms to human preferences, not by how human-AI interaction evolves. While RLHF produces models that generate high-quality individual responses, it does not account for conversation dynamics---how users learn AI capabilities over turns or how AI systems respond to evolving query patterns.

This work complements RLHF by proposing conversation-level metrics that measure alignment during interaction. Rather than ''does the AI produce responses humans prefer?'', the question becomes ''do users and AI co-adapt during conversations?''

\subsubsection{User Adaptation and Learning Curves}

Human-computer interaction research recognizes that users learn to communicate more effectively with systems over time \citep{newell1981mechanisms}. Learning curve theory predicts that task performance improves with practice. In conversational AI contexts, this manifests as query reformulation patterns: novice users reformulate frequently while experienced users craft effective queries initially.

Prior work on query reformulation has focused on search engines and information retrieval systems \citep{huang2009query}. Users reformulate queries when initial results are unsatisfactory. However, this literature treats the search engine as static; reformulation reflects user learning about the corpus and ranking algorithm, not about an adaptive AI agent.

More recent work examines how users adapt to AI assistants. Studies of voice interfaces \citep{luger2016like} find that users develop mental models of system capabilities through trial-and-error. Conversational repair strategies \citep{ashktorab2019resilient} show users employ systematic patterns when AI misunderstands them. These findings suggest users actively learn AI capabilities during conversations, but existing work has not quantified learning rates or connected them to alignment quality.

This contribution operationalizes user learning through reformulation slope---the rate at which reformulation frequency decreases over conversation turns. Negative slope indicates learning, while flat or positive slope suggests no learning or engagement decay.

\subsubsection{Conversational AI Evaluation}

Evaluation of multi-turn conversational systems has traditionally focused on task completion, dialogue coherence, and user satisfaction \citep{walker1997paradise}. Task-oriented dialogue systems are evaluated on success rate and efficiency. Open-domain chatbots are evaluated on engagement, appropriateness, and subjective quality ratings.

Recent work has introduced conversation-level metrics beyond aggregate ratings. Mehri \& Eskenazi (2020) propose USR, an automatic metric that correlates with human judgments of naturalness and coherence. Zhang et al. (2021) analyze turn-level topic shifts and introduce metrics for conversation depth and breadth. Durmus et al. (2023) evaluate helpfulness in multi-turn contexts.

However, these metrics still treat the AI as the primary unit of analysis---measuring response quality, coherence, or task enablement. They do not measure user behavioral changes during conversation or coupling between user and AI adaptation patterns.

Interactive Alignment Theory from psycholinguistics \citep{pickering2004toward} provides a relevant framework. In human-human dialogue, both participants unconsciously align their linguistic representations through priming and adaptation. Successful conversations exhibit high alignment on multiple levels. While this theory describes human cognition, the core insight applies to human-AI interaction: alignment emerges through mutual adaptation during conversation.

This work extends conversation evaluation by introducing bidirectional metrics. It measures user learning (reformulation slope) alongside AI responsiveness (diversity correlation) and proposes coupling strength (correlation between these signals) as a conversation-level alignment metric.

\subsubsection{Positioning This Contribution}

No prior work has operationalized bidirectional alignment through observable conversation behaviors. RLHF measures alignment as policy optimization against human preferences. HCI learning curve research measures user adaptation but treats AI as static. Conversation evaluation metrics measure output quality but not behavioral dynamics. This work combines insights from all three areas to introduce behavioral coupling as a metric that quantifies co-adaptation between users and AI during conversations.
